% 00-about.tex: что представляет собой настоящий документ и как читать его метки.
\section*{О настоящей версии}
\addcontentsline{toc}{section}{О настоящей версии}

Версия 1 настоящего отчёта (\code{полярис\_отчёт\_о\_проекте.пдф}, сохранена без изменений рядом с
этим файлом) описывала Полярис как проект базы данных: двенадцать таблиц вокруг одного
удостоверения, модель <<сущность--связь>>, автомат жизненного цикла и довод о том, что
идентификационные данные обязаны быть постквантовыми с первого дня. Версия 2
(\code{полярис\_отчёт\_о\_проекте\_в2.пдф}, также сохранена) описывала систему в том виде, в каком
она существовала на теге \code{в9.345}: тридцать шесть таблиц, протокол федерации, слой
прозрачности и ткань межведомственного обмена; она была пересмотрена на \code{в9.355}, когда
работа на стороне держателя, отмеченная в ней как отсутствующая, была выполнена. Версия 3
описывает систему на версии \code{1.0.0-кв.7} (коммит \code{4e74376}, 20 сентября 2026 года),
измеренную заново целиком, а не залатанную, и написана она для читателя, который никогда не
открывал репозиторий. Впервые она была написана на \code{1.0.0-кв.1} четырьмя днями ранее и
перемерена на месте; четырнадцать из двадцати шести величин в приложении~\ref{app:counts} не
сдвинулись, а те, что сдвинулись, там отмечены.

\textbf{О передаче обозначений.} Настоящее издание написано для читателя, не владеющего английским языком, и потому по-русски передан не только текст, но и каждое обозначение: имена таблиц, проверок, полей, путей, пакетов и стандартов. В самом репозитории эти обозначения записаны латиницей, и тот, кто пойдёт от этого документа к коду, найдёт их там под латинскими именами; полное соответствие хранится в словаре рядом с исходным текстом документа, а сокращения раскрыты в приложении~\ref{app:glossary}. Без перевода оставлены две вещи: девиз проекта, написанный по-латыни, и номер коммита, который есть шестнадцатеричное число, а не слово.

\textbf{Что изменилось, в одном абзаце.} Между версией 2 и настоящей репозиторий перестал
считать собственный прогресс и начал считать свидетелей. Четыре самостоятельных пакета, которые
устанавливает доверяющая сторона, находятся в публичных реестрах. Кошелёк, написанный
посторонними людьми и запущенный без изменений, предъявил удостоверение верификатору и был
принят, после того как один раз получил отказ из-за дефекта, мимо которого прошёл каждый
внутренний инструмент. Размещённый пакет проверок соответствия Фонда ОпенАйДи прогнал профиль
верификатора через открытый интернет. Сам слой проверки подвергся атаке изнутри: инвариантные
проверки, ограничения, триггеры, уникальные индексы, хранимые процедуры, свидетельства с нулевым
разглашением, контракт соответствия и оба эталонных верификатора были поочерёдно мутированы и
обязаны были покраснеть, и многие поначалу этого не сделали. Физический слой стал предметом,
который способен изготовить кто-то другой. И два слова на парадной двери были ослаблены:
\emph{постквантовый} до \emph{алгоритмической гибкости при проверяемом пути миграции}, а
\emph{устойчивый к принуждению} до \emph{осведомлённого о принуждении}, потому что имеющиеся
свидетельства не подтверждали более сильные формулировки. Приложение~\ref{app:from-v2}
перечисляет каждое утверждение версии 2, которое настоящая версия отзывает или сужает.

\textbf{Три даты, названные прямо.} Основной текст версии 2 был измерен на \code{в9.345}; её
разделы о стороне держателя на \code{в9.355}. Всё в настоящей версии, каждая величина, каждая
метка и каждый рисунок, измерено на \code{1.0.0-кв.7} методами из приложения~\ref{app:counts}.
Там, где число старше этого, потому что репозиторий не перезапускал породивший его инструмент,
документ говорит об этом рядом с числом. Ничто из работавшего на \code{в9.355} работать не
перестало.

Документ сохраняет то, что версия 1 угадала верно: геометрию схемы с токеном в центре,
используемую здесь как центр каждой карты, и автомат жизненного цикла; а также то, что верно
угадала версия 2: концентрическое прочтение системы, по одному кольцу на раздел, где каждое
кольцо вводится отказом кольца внутри него.

\subsection*{Пять меток}

Каждая способность несёт одну из пяти меток, чтобы работающий код никогда не путали с
намерением, а то, что проверил автор, никогда не путали с тем, что проверил кто-то ещё.

\begin{description}
  \item[\sImpl] Путь, который выполняется на \code{1.0.0-кв.7}, задействован тестом или учением
    в непрерывной интеграции и закреплён инвариантной проверкой, роняющей сборку, если он
    перестаёт быть истинным.
  \item[\sExp] Путь, который выполняется, но ограничен способом, который называет сам
    репозиторий: не покрыт НИ, условен по масштабу или зависит от серверной части, которую НИ
    имитирует.
  \item[\sExt] Свойство, которое репозиторий реализует и проверяет внутренними средствами, но
    которое никакая внешняя сторона не проверяла аудитом, не атаковала, не разворачивала и не
    реализовывала независимо заново.
  \item[\sWit] Путь, который названная сторона за пределами репозитория действительно
    задействовала, в конкретную дату, с результатом, на который можно указать: кошелёк, служба
    проверки соответствия, опубликованный набор тестов. Метка намеренно редка, она даётся только
    там, где в собственной сводной таблице репозитория заполнена строка, и она никогда не
    означает аудит, сертификацию или пилотное внедрение. В настоящем документе её несут три
    вещи.
  \item[\sFut] Предложение: строка дорожной карты, проектная записка или замысел, существующий
    только в настоящем документе. Это не код.
\end{description}

Строка, несущая две метки, работает сегодня и всё же требует внешней проверки для более
сильного утверждения либо была задействована извне на одном пути и не была на остальных.

Числа измерены на дереве на \code{1.0.0-кв.7} методами из приложения~\ref{app:counts}, а не
переписаны из более ранних документов. Там, где собственные документы репозитория расходятся с
его кодом, побеждает код, а расхождение фиксируется в авторских заметках, сопровождающих
настоящую версию.

\subsection*{Аналогия и её предел}

Документ объясняет Полярис через одну сквозную аналогию: город, обнесённый стеной. Запись
удостоверения есть реестр городской ратуши; ограничения есть стена; конечные точки есть ворота;
журналы прозрачности есть сторожевые башни; независимые свидетели стоят за стеной и наблюдают за
башнями; другие удостоверяющие органы есть другие города, соединённые дорогами с односторонним
движением; над всем городом стоят законы, связывающие и сам город; и, впервые в настоящей
версии, к воротам пришли путники издалека, и найденное ими есть единственный отчёт, который
город написал не о себе сам. Аналогия описывает расположение и линии обзора, а не политическую
централизацию: Полярис федеративен в пределах каждого удостоверяющего органа, без центральной
базы данных, корня или службы проверки, и всякое возвращение аналогии возвращается к
действительному механизму под его настоящим именем.

\subsection*{Один тезис и одно правило}

Система есть попытка построить идентификационную инфраструктуру, в которой важным утверждениям
не приходится доверять лишь потому, что центральное приложение называет их истинными.
Полномочия, состояние, история, границы приватности и собственное поведение системы
последовательно выставляются на независимую проверку и стесняются слоями, которые не разделяют
допущений друг друга. Каждый слой существует потому, что слой внутри него может быть совершенно
верен и всё же не доказать того, что обществу необходимо знать. Эта последовательность есть
хребет документа.

Правило, принятое 13 сентября 2026 года и применяемое ко всему с тех пор, состоит в том, что
единицей прогресса является новая пережитая внешняя зависимость: названный кошелёк, названный
профиль соответствия, названная доверяющая сторона, дефект, заведённый кем-то, кто не является
автором. Число инвариантов, версий и строк прогрессом не является, и сводная таблица репозитория
отказывается его записывать. Настоящий документ следует этому правилу. Там, где он сообщает
внутренние свидетельства, он говорит об этом и держит эти свидетельства вне предложений,
сообщающих свидетельства другого рода.
